\section{Methodology}
\label{sec:method}

\subsection{Research Questions and Hypotheses}

Our main question is: \emph{Do multi-agent LLM systems with capability-level diversity
(heterogeneous model sizes within a single family) demonstrate superior robustness
on adversarially-designed reasoning tasks compared to homogeneous ensembles,
due to reduced correlated error rates?}

We decompose this into four testable hypotheses:

\begin{itemize}[leftmargin=*,itemsep=2pt,topsep=2pt]
    \item \Hone (\emph{Accuracy}): Heterogeneous ensembles achieve higher accuracy
          than homogeneous ensembles.
    \item \Htwo (\emph{Error Correlation}): Error correlation between
          \haiku--\haiku pairs exceeds that between \haiku--\sonnet pairs.
    \item \Hthree (\emph{Debate Benefit}): Structured debate between agents of
          different capability levels improves accuracy over voting alone.
    \item \Hfour (\emph{Adversarial Robustness}): Heterogeneous ensembles degrade
          less under paraphrase perturbations than homogeneous ensembles.
\end{itemize}

\subsection{Adversarial Benchmark: \advbench}

We construct a 30-question adversarial reasoning benchmark
spanning five categories of six questions each.
Each category is designed to expose a specific cognitive failure mode.

\para{Category 1: Misleading Math.}
Problems where a plausible arithmetic shortcut gives the wrong answer.
The canonical example is the bat-and-ball problem: a bat and ball together cost
\$1.10, and the bat costs exactly \$1.00 more---the na\"ive answer is 10 cents,
but the correct answer is 5 cents. We include six such problems exploiting
rate calculations, percentage comparisons, and compound interest.

\para{Category 2: Causal Traps.}
Problems involving correlation-causation confusion, hidden confounders,
and Simpson's paradox variants.
For example: Hospital A has a 2\% death rate and Hospital B has a 1\% death rate---should
one choose Hospital B? The correct answer is no, since Hospital A may treat sicker patients
(Simpson's paradox).

\para{Category 3: Logical Deception.}
Multi-step logic puzzles with misleading surface patterns, inspired by
Knight-Knave-style problems~\citep{wu2025debate}.
Surface cues suggest one answer, but the full logical deduction yields another.
We include a light-switch heat detection puzzle (which bulb was recently on?)
and a classic rope-burning timing puzzle.

\para{Category 4: Numerical Tricks.}
Problems exploiting anchoring, base rate neglect, and unit confusion.
The canonical example is the medical test puzzle: if 1\% of people have disease~X
and a test is 99\% accurate, a positive result implies only $\approx$50\% probability
of having the disease (by Bayes' theorem), not 99\% as anchoring suggests.

\para{Category 5: Framing Effects.}
Equivalent problems presented under different framings to test reasoning consistency.
The key example is surgery recommendation: ``A surgery has a 90\% survival rate---would
you recommend it?'' The correct reasoning requires clinical context, but different
framings should not change the reasoning process.

\para{Paraphrase variants.}
We additionally create six paraphrase variants of selected questions
to test robustness under surface-level rewording.
The same logical structure is preserved, but vocabulary and sentence structure differ.

\para{Benchmark validity.}
All 30 questions have verified ground-truth answers derived from first principles.
Two questions (light-switch heat detection and rope-burning timing) require procedural
free-text answers and are evaluated using content-based matching.
One question (exponential population doubling) contains a genuine linguistic ambiguity
in the word ``exceed''; we discuss this as a benchmark design lesson in \secref{sec:results}.

\subsection{Experimental Conditions}

We evaluate five conditions representing a diversity continuum
(\tabref{tab:conditions}):

\begin{table}[t]
\centering
\caption{Five experimental conditions. ``H'' denotes \haiku; ``S'' denotes \sonnet.
Aggregation applies to 3-agent conditions only.}
\label{tab:conditions}
\resizebox{\textwidth}{!}{%
\begin{tabular}{@{}llccc@{}}
\toprule
Condition & Description & Agents & Models & Aggregation \\
\midrule
\singlehaiku   & Single small model           & 1 & 1H      & Direct answer \\
\singlesonnet  & Single large model           & 1 & 1S      & Direct answer \\
\homoghaiku    & Homogeneous small ensemble   & 3 & 3H      & Majority vote \\
\hetero        & Heterogeneous ensemble       & 3 & 2H + 1S & Majority vote \\
\debate        & Structured debate            & 2 & 1H + 1S & H revises after S critiques \\
\bottomrule
\end{tabular}
}
\end{table}

The \hetero condition uses two \haiku agents and one \sonnet agent, so \haiku
has a majority position by count. This tests whether the stronger model's
minority opinion can override the weaker majority, or whether majority voting
conservatively favors the smaller model.

The \debate condition implements a three-turn protocol:
(1) \haiku proposes an answer with chain-of-thought reasoning;
(2) \sonnet critiques the proposal using a dedicated critic system prompt;
(3) \haiku revises its answer with a dedicated reviser system prompt.
The final answer is \haiku's revised response.

\subsection{Models and Implementation}

We use the Anthropic API with two models:
\haiku (\texttt{claude-haiku-4-5-20251001}), a fast cost-efficient model,
and \sonnet (\texttt{claude-sonnet-4-6}), a more capable model.

All API calls use temperature 0.0 for deterministic, reproducible results,
with a maximum of 512 tokens per response.
The system prompt for all non-debate conditions encourages step-by-step reasoning
and instructs agents to state \texttt{FINAL ANSWER:} clearly at the end.
The debate uses separate system prompts for the critic and reviser roles.
We set the random seed to 42 for all Python operations,
and include a 0.3--0.5 second inter-call delay for rate limit compliance.
Failed API calls are retried up to three times with exponential backoff.

\subsection{Evaluation Metrics}

\para{Task accuracy.}
The primary metric is the fraction of questions answered correctly.
Ground truth is a verified scalar answer for most questions and
a content-based match for procedural questions.
We report 95\% bootstrap confidence intervals using 1000 resamples.

\para{Error correlation.}
For each pair of conditions $A$ and $B$, we compute the Pearson correlation
between their binary error vectors $\ve_A, \ve_B \in \{0,1\}^{30}$,
where $e_i = 1$ if condition $i$ answers question $j$ incorrectly.
Lower correlation indicates more diverse failure modes.
We report the diversity index as $1 - r$.

\para{Debate improvement rate.}
The fraction of questions where the debate protocol changes \haiku's answer
from wrong to right (improvement), right to wrong (degradation), or both.

\para{Adversarial robustness.}
The accuracy drop between original and paraphrase variants (6 questions).
Smaller drops indicate higher robustness under surface perturbation.

\para{Statistical tests.}
We use McNemar's test for pairwise accuracy comparisons between conditions
and report two-tailed $p$-values.
Pearson correlation significance is assessed via $t$-test with $n{-}2$ degrees of freedom.
